\section{Модели машинного обучения для декодирования EEG}

После этапов предобработки и извлечения признаков задача декодирования EEG сводится к построению отображения между многоканальным временным сигналом и целевой переменной. В контексте интерфейсов мозг-компьютер (brain-computer interface, BCI) эта переменная часто имеет вид класса, однако те же принципы применимы и к задачам регрессии или многовыходного предсказания. Для настоящей работы особенно важны модели, устойчивые к малому числу обучающих примеров, шумности сигнала и межсубъектной вариативности.

\subsection{Классические модели машинного обучения}

В EEG-BCI долгое время широко применялись линейный дискриминантный анализ (linear discriminant analysis, LDA), логистическая регрессия (logistic regression), метод опорных векторов (support vector machine, SVM), случайный лес (random forest) и их регуляризованные варианты. Их преимущество заключается в умеренной вычислительной сложности, сравнительно высокой интерпретируемости и возможности работы с заранее заданными признаками: band power, ковариационными матрицами, признаками CSP, временно-частотными характеристиками.

Обзор алгоритмов классификации EEG-BCI показывает, что при малых обучающих выборках особенно важны регуляризация, устойчивые оценки ковариаций, адаптация и transfer learning~\footcite{10.1088/1741-2552/aab2f2}. Это соответствует условиям задач когнитивного декодирования, где число повторений обычно ограничено, а индивидуальные различия между испытуемыми могут быть сопоставимы по величине с целевыми эффектами.

Для реконструкции бинарного изображения $6 \times 6$ такие методы могут использоваться как базовые модели: например, отдельный классификатор или регрессор может предсказывать состояние каждого пикселя, а затем результаты объединяются в двумерную структуру. Достоинство такого подхода состоит в контролируемой сложности модели и возможности анализировать вклад пространственно-спектральных признаков в отдельные элементы изображения.

\subsection{Нейросетевые модели и EEGNet}

Глубокие нейронные сети позволяют обучать представления EEG непосредственно из данных или из слабо обработанных временных окон. Свёрточные нейронные сети (convolutional neural networks, CNN) применяются для выделения локальных временных, частотных и пространственных паттернов. Компактные архитектуры семейства EEGNet ориентированы на EEG-задачи и стремятся совместить временные фильтры, пространственную фильтрацию по каналам и классификацию в единой end-to-end модели.

Преимущество таких моделей заключается в способности автоматически подбирать признаки, однако это преимущество проявляется только при достаточном объёме и разнообразии данных. В задачах с малым числом эпох глубокие модели могут переобучаться на артефакты, индивидуальные особенности испытуемых или особенности конкретной сессии. Поэтому в практических BCI-сценариях нейросетевые модели часто дополняются временной агрегацией предсказаний, скользящими окнами и специальными схемами валидации~\footcite{10.1109/jsen.2023.3270281}.

Для настоящей работы CNN- и EEGNet-подобные архитектуры являются важным ориентиром, но не снимают необходимости в простых базовых моделях. При реконструкции визуального образа из EEG, зарегистрированной в фазе воспоминания, целевой сигнал слабее и менее стабилен, чем при задачах моторного воображения или вызванных потенциалов, поэтому рост сложности модели должен быть обоснован результатами валидации.

\subsection{Transformer-based models и foundation models для EEG}

Развитие трансформерных архитектур и самообучения (self-supervised learning) привело к появлению EEG foundation models, обучаемых на больших массивах разнородных EEG-данных. В отличие от моделей, обучаемых только на одном датасете и одной задаче, такие модели стремятся сформировать универсальные представления EEG, которые затем адаптируются к downstream-задачам. Критический обзор ранних EEG foundation models подчёркивает, что ключевыми вопросами остаются представление входного сигнала, выбор self-supervised objective и реалистичная стратегия оценки переносимости~\footcite{10.1088/1741-2552/ae4455}.

Модель LaBraM (Large Brain Model) предлагает представлять EEG через channel patches и обучать трансформер на восстановлении маскированных нейронных кодов. Авторы позиционируют её как модель для cross-dataset learning и обучают на примерно 2500 часах EEG из разных наборов данных~\footcite{10.48550/arXiv.2405.18765}. Neuro-GPT сочетает EEG-энкодер и GPT-подобную модель, используя self-supervised восстановление маскированных EEG-сегментов, а затем дообучается на задаче классификации моторного воображения в условиях малого числа испытуемых~\footcite{10.48550/arXiv.2311.03764}.

CBraMod развивает это направление за счёт criss-cross transformer architecture, разделяющей пространственное и временное моделирование EEG, и проверяется на нескольких downstream BCI-задачах~\footcite{10.48550/arXiv.2412.07236}. REVE делает акцент на переносимости между различными схемами регистрации: модель использует 4D positional encoding и masked autoencoding objective, что позволяет обрабатывать сигналы с разной длиной и различным расположением электродов~\footcite{10.48550/arXiv.2510.21585}.

Для данной магистерской работы эти модели важны прежде всего как современный контекст развития EEG-декодирования. Они показывают общий тренд к масштабному предобучению и переносимым представлениям, но не являются прямым решением рассматриваемой задачи. Во-первых, их обучение требует крупных корпусов EEG, значительно превосходящих объём данных типичного эксперимента по визуальному воображению. Во-вторых, большинство демонстраций относится к классификации или клиническим задачам, а не к пиксельной реконструкции малых визуальных стимулов из состояния воспоминания. Поэтому foundation models следует рассматривать как перспективное направление для будущего расширения, тогда как в основной постановке оправдано использование контролируемых по сложности моделей и признаков.

\subsection{Стратегии обучения: intra-subject и inter-subject}

Ключевым методологическим вопросом в EEG-декодировании является выбор стратегии обучения и оценки. Внутри-субъектное обучение (intra-subject или within-subject learning) предполагает, что модель обучается и тестируется на данных одного испытуемого при корректном разделении эпох на обучающую и тестовую части. Такой сценарий обычно даёт более высокое качество, поскольку модель может адаптироваться к индивидуальной топографии, амплитудам и спектральным особенностям EEG.

Межсубъектное обучение (cross-subject learning) проверяет переносимость модели между людьми: часть испытуемых используется для обучения, а другой испытуемый или группа испытуемых оставляется для тестирования. Этот сценарий ближе к практическому применению, поскольку снижает потребность в длительной индивидуальной калибровке, но обычно сложнее из-за различий в анатомии, расположении электродов, уровне шума и когнитивных стратегиях.

Для задач реконструкции визуальных стимулов различие между этими стратегиями особенно важно. Высокое качество внутри одного испытуемого может отражать не только устойчивое нейронное представление стимула, но и индивидуальные особенности данных. Напротив, межсубъектная оценка демонстрирует, насколько модель улавливает общие закономерности, переносимые между людьми. Поэтому при интерпретации результатов необходимо явно указывать, какой протокол использован, и не переносить выводы из intra-subject сценария на inter-subject применение без дополнительной проверки.

\subsection*{Выводы по разделу}

Модели машинного обучения для EEG-декодирования образуют спектр от простых интерпретируемых классификаторов до крупных foundation models. Классические методы остаются значимыми благодаря устойчивости при малых выборках и возможности работы с пространственно-спектральными признаками. Нейросетевые и трансформерные подходы расширяют возможности автоматического обучения представлений, однако требуют строгого контроля переобучения и реалистичной оценки переносимости. В контексте реконструкции бинарных изображений из EEG visual imagery наиболее обоснованной является стратегия, в которой сложность модели соотносится с объёмом данных, а результаты анализируются отдельно для intra-subject и inter-subject протоколов.
